\documentclass[10pt,a4paper]{amsart}
\usepackage{amsmath,amssymb,amsthm,graphicx}
% Retained preamble for the Stacks Project chapter {CH} (shared project preamble).
% The chapter file pinned-original-source.tex only carries \input{{preamble}} on its first line,
% so no preamble text lives inside the chapter itself.  The Stacks Project uses one single shared
% preamble for every chapter: the snapshot stages it three times (sources/stacks-chow/preamble.tex,
% sources/stacks-etale-cohomology/preamble.tex and focused-seeds/AGP/sources/stacks-snapshot-20261011/
% preamble.tex) and those copies are identical in size.  The block below is copied mechanically,
% line for line, from sources/stacks-chow/preamble.tex lines 14-18 and 154-217: the Xy-pic setup and
% the native theorem declarations together with the mathematical macros.  Nothing was added, removed
% or redefined; the wrapper class and the amsmath/amssymb/amsthm/graphicx base packages are supplied
% by the assembler.
% For commutative diagrams we use Xy-pic
% Chapter: stacks-spaces-more-morphisms.tex -- More on Morphisms of Spaces.
\usepackage[all]{xy}

% We use 2cell for 2-commutative diagrams.
\xyoption{2cell}

% Theorem environments.
%
\theoremstyle{plain}
\newtheorem{theorem}[subsection]{Theorem}
\newtheorem{proposition}[subsection]{Proposition}
\newtheorem{lemma}[subsection]{Lemma}

\theoremstyle{definition}
\newtheorem{definition}[subsection]{Definition}
\newtheorem{example}[subsection]{Example}
\newtheorem{exercise}[subsection]{Exercise}
\newtheorem{situation}[subsection]{Situation}

\theoremstyle{remark}
\newtheorem{remark}[subsection]{Remark}
\newtheorem{remarks}[subsection]{Remarks}

\numberwithin{equation}{subsection}

% Macros
%
\def\lim{\mathop{\mathrm{lim}}\nolimits}
\def\colim{\mathop{\mathrm{colim}}\nolimits}
\def\Spec{\mathop{\mathrm{Spec}}}
\def\Hom{\mathop{\mathrm{Hom}}\nolimits}
\def\Ext{\mathop{\mathrm{Ext}}\nolimits}
\def\SheafHom{\mathop{\mathcal{H}\!\mathit{om}}\nolimits}
\def\SheafExt{\mathop{\mathcal{E}\!\mathit{xt}}\nolimits}
\def\Sch{\mathit{Sch}}
\def\Mor{\mathop{\mathrm{Mor}}\nolimits}
\def\Ob{\mathop{\mathrm{Ob}}\nolimits}
\def\Sh{\mathop{\mathit{Sh}}\nolimits}
\def\NL{\mathop{N\!L}\nolimits}
\def\CH{\mathop{\mathrm{CH}}\nolimits}
\def\proetale{{pro\text{-}\acute{e}tale}}
\def\etale{{\acute{e}tale}}
\def\QCoh{\mathit{QCoh}}
\def\Ker{\mathop{\mathrm{Ker}}}
\def\Im{\mathop{\mathrm{Im}}}
\def\Coker{\mathop{\mathrm{Coker}}}
\def\Coim{\mathop{\mathrm{Coim}}}

% Boxtimes
%
\DeclareMathSymbol{\boxtimes}{\mathbin}{AMSa}{"02}

%
% Macros for moduli stacks/spaces
%
\def\QCohstack{\mathcal{QC}\!\mathit{oh}}
\def\Cohstack{\mathcal{C}\!\mathit{oh}}
\def\Spacesstack{\mathcal{S}\!\mathit{paces}}
\def\Quotfunctor{\mathrm{Quot}}
\def\Hilbfunctor{\mathrm{Hilb}}
\def\Curvesstack{\mathcal{C}\!\mathit{urves}}
\def\Polarizedstack{\mathcal{P}\!\mathit{olarized}}
\def\Complexesstack{\mathcal{C}\!\mathit{omplexes}}
% \Pic is the operator that assigns to X its picard group, usage \Pic(X)
% \Picardstack_{X/B} denotes the Picard stack of X over B
% \Picardfunctor_{X/B} denotes the Picard functor of X over B
\def\Pic{\mathop{\mathrm{Pic}}\nolimits}
\def\Picardstack{\mathcal{P}\!\mathit{ic}}
\def\Picardfunctor{\mathrm{Pic}}
\def\Deformationcategory{\mathcal{D}\!\mathit{ef}}
\usepackage{cleveref}
\crefname{originallemma}{original lemma}{original lemmas}
\Crefname{originallemma}{Original lemma}{Original lemmas}
\makeatletter
\newlabel{lemma-chow-noetherian-separated}{{088U}{source}{}{sourcepointer-1}{}}
\newlabel{lemma-chow-noetherian-separated@cref}{{[originallemma][0][]088U}{[1][0][]source}{}{}{}}
\makeatother
\crefname{originalreference}{original reference}{original references}
\Crefname{originalreference}{Original reference}{Original references}
\makeatletter
\newlabel{spaces-limits-lemma-descend-finite-presentation}{{07SK}{source}{}{sourcepointer-2}{}}
\newlabel{spaces-limits-lemma-descend-finite-presentation@cref}{{[originalreference][0][]07SK}{[1][0][]source}{}{}{}}
\makeatother
\crefname{originalreference}{original reference}{original references}
\Crefname{originalreference}{Original reference}{Original references}
\makeatletter
\newlabel{spaces-limits-lemma-descend-separated}{{084T}{source}{}{sourcepointer-3}{}}
\newlabel{spaces-limits-lemma-descend-separated@cref}{{[originalreference][0][]084T}{[1][0][]source}{}{}{}}
\makeatother
\crefname{originalreference}{original reference}{original references}
\Crefname{originalreference}{Original reference}{Original references}
\makeatletter
\newlabel{spaces-limits-lemma-descend-separated-morphism}{{0851}{source}{}{sourcepointer-4}{}}
\newlabel{spaces-limits-lemma-descend-separated-morphism@cref}{{[originalreference][0][]0851}{[1][0][]source}{}{}{}}
\makeatother
\crefname{originalreference}{original reference}{original references}
\Crefname{originalreference}{Original reference}{Original references}
\makeatletter
\newlabel{spaces-limits-proposition-approximate}{{07SU}{source}{}{sourcepointer-5}{}}
\newlabel{spaces-limits-proposition-approximate@cref}{{[originalreference][0][]07SU}{[1][0][]source}{}{}{}}
\makeatother
\crefname{originalreference}{original reference}{original references}
\Crefname{originalreference}{Original reference}{Original references}
\makeatletter
\newlabel{spaces-limits-proposition-separated-closed-in-finite-presentation}{{0873}{source}{}{sourcepointer-6}{}}
\newlabel{spaces-limits-proposition-separated-closed-in-finite-presentation@cref}{{[originalreference][0][]0873}{[1][0][]source}{}{}{}}
\makeatother
\title{Chow's lemma with an immersion into projective space for separated finite type morphisms of algebraic spaces}
\author{The Stacks Project Authors}
\date{stacks-project:spaces-more-morphisms.tex:local-snapshot-20261011; GNU FDL 1.2 or later}
\begin{document}
\maketitle
\noindent\textit{Editorial note.} Counted span is the author's own proof environment, source lines 9446-9477, immediately following the statement at 9431-9444.\par\medskip
\noindent\textit{Editorial note.} Two 'Clearly' transfer steps (source lines 9450-9452 and 9471-9473) are retained verbatim; the justification recorded in proof\_roadmap (composition preserves proper surjective, base change preserves immersions) is the reviewer's checklist, not author text.\par\medskip
\noindent\textit{Editorial note.} All five cross-chapter pointers are Limits of Spaces lemmas recorded in assembly.external\_references; the terminal same-chapter pointer is a line-indexed stub into the pinned source.\par\medskip
\noindent\textit{Allowed standard prerequisites (editorial).}\begin{itemize}
\item Finite presentation replacement by a closed immersion. Limits of Spaces, Proposition spaces-limits-proposition-separated-closed-in-finite-presentation, cited at proof lines 9448-9450. X embeds as a closed subspace of X\_1 separated and of finite presentation over Y; proper surjective X' -> X composed with the closed immersion transfers the conclusion back.
\item Approximation of separated quasi-compact spaces by finite type layers. Limits of Spaces, Proposition spaces-limits-proposition-approximate, cited at proof lines 9455-9457. Writes Y = lim Y\_i with Y\_i quasi-separated of finite type over Spec(Z), the ladder on which everything descends.
\item Descending separatedness and finite presentation along the limit. Limits of Spaces, Lemmas spaces-limits-lemma-descend-separated (9457-9459), spaces-limits-lemma-descend-finite-presentation (9462-9465), spaces-limits-lemma-descend-separated-morphism (9468-9470). Produces an index i and X\_i -> Y\_i of finite presentation, separated, with X = Y x\_\{Y\_i\} X\_i.
\end{itemize}\medskip
\ifdefined\hypertarget\hypertarget{sourcepointer-1}{}\fi
\noindent\textit{Source reference.} Original lemma 088U: sources/stacks-spaces-more-morphisms/pinned-original-source.tex, source lines 9275--9275. Navigation pointer only: the counted proof cites this statement by \textbackslash{}ref as a standing input; the deduction the counted proof has to make is carried out inside the counted span, and no part of that deduction is delegated to this pointer.\par\medskip
\ifdefined\hypertarget\hypertarget{sourcepointer-2}{}\fi
\noindent\textit{Source reference.} Original reference 07SK: focused-seeds/geometry/S9-spaces-chow-finite-type-separated/ref-tags.txt, source lines 5--5. Navigation pointer only: the counted proof cites this statement by \textbackslash{}ref as a standing input; the deduction the counted proof has to make is carried out inside the counted span, and no part of that deduction is delegated to this pointer.\par\medskip
\ifdefined\hypertarget\hypertarget{sourcepointer-3}{}\fi
\noindent\textit{Source reference.} Original reference 084T: focused-seeds/geometry/S9-spaces-chow-finite-type-separated/ref-tags.txt, source lines 6--6. Navigation pointer only: the counted proof cites this statement by \textbackslash{}ref as a standing input; the deduction the counted proof has to make is carried out inside the counted span, and no part of that deduction is delegated to this pointer.\par\medskip
\ifdefined\hypertarget\hypertarget{sourcepointer-4}{}\fi
\noindent\textit{Source reference.} Original reference 0851: focused-seeds/geometry/S9-spaces-chow-finite-type-separated/ref-tags.txt, source lines 7--7. Navigation pointer only: the counted proof cites this statement by \textbackslash{}ref as a standing input; the deduction the counted proof has to make is carried out inside the counted span, and no part of that deduction is delegated to this pointer.\par\medskip
\ifdefined\hypertarget\hypertarget{sourcepointer-5}{}\fi
\noindent\textit{Source reference.} Original reference 07SU: focused-seeds/geometry/S9-spaces-chow-finite-type-separated/ref-tags.txt, source lines 8--8. Navigation pointer only: the counted proof cites this statement by \textbackslash{}ref as a standing input; the deduction the counted proof has to make is carried out inside the counted span, and no part of that deduction is delegated to this pointer.\par\medskip
\ifdefined\hypertarget\hypertarget{sourcepointer-6}{}\fi
\noindent\textit{Source reference.} Original reference 0873: focused-seeds/geometry/S9-spaces-chow-finite-type-separated/ref-tags.txt, source lines 9--9. Navigation pointer only: the counted proof cites this statement by \textbackslash{}ref as a standing input; the deduction the counted proof has to make is carried out inside the counted span, and no part of that deduction is delegated to this pointer.\par\medskip
\begin{lemma}
\label{lemma-chow-finite-type-separated}
Let $S$ be a scheme. Let $f : X \to Y$ be a morphism of algebraic spaces
over $S$. Assume $f$ separated of finite type, and $Y$ separated and
quasi-compact. Then there exists a commutative diagram
$$
\xymatrix{
X \ar[rd] & X' \ar[l] \ar[d] \ar[r] & \mathbf{P}^n_Y \ar[ld] \\
& Y
}
$$
where $X' \to X$ is proper surjective morphism
and the morphism $X' \to \mathbf{P}^n_Y$ is an immersion.
\end{lemma}

\begin{proof}
By
Limits of Spaces, Proposition
\ref{spaces-limits-proposition-separated-closed-in-finite-presentation}
we can find a closed immersion $X \to X_1$ where $X_1$ is separated
and of finite presentation over $Y$. Clearly, if we prove the assertion
for $X_1 \to Y$, then the result follows for $X$. Hence we may assume that
$X$ is of finite presentation over $Y$.

\medskip\noindent
We may and do replace the base scheme $S$ by $\Spec(\mathbf{Z})$.
Write $Y = \lim_i Y_i$ as a directed limit of
quasi-separated algebraic spaces of finite type over $\Spec(\mathbf{Z})$, see
Limits of Spaces,
Proposition \ref{spaces-limits-proposition-approximate}.
By Limits of Spaces, Lemma \ref{spaces-limits-lemma-descend-separated}
we may assume that $Y_i$ is separated for all $i$.
By
Limits of Spaces,
Lemma \ref{spaces-limits-lemma-descend-finite-presentation}
we can
find an index $i \in I$ and a scheme $X_i \to Y_i$ of finite presentation
so that $X = Y \times_{Y_i} X_i$.
By
Limits of Spaces,
Lemma \ref{spaces-limits-lemma-descend-separated-morphism}
we may assume that $X_i \to Y_i$ is separated.
Clearly, if we prove the assertion for
$X_i$ over $Y_i$, then the assertion holds for $X$. The case
$X_i \to Y_i$ is treated by
Lemma \ref{lemma-chow-noetherian-separated}.
\end{proof}

\end{document}
